% ============================================================
%  Modular arithmetic: the "numbers" we compute with
% ============================================================

\section{Modular arithmetic and finite fields}

All the computations on the curve are done \textbf{modulo a prime} $p$: we never work with real numbers, only with the integers $\{0, 1, \ldots, p-1\}$, and every result is reduced by taking the remainder of the division by $p$.

\begin{frameddefn}[label={def:congruence}]{Congruence modulo $p$}
  Two integers $a, b$ are \textit{congruent modulo} $p$, written $a \equiv b \pmod p$, if $p$ divides $a - b$. Every integer is congruent to exactly one element of $\{0, \ldots, p-1\}$, its remainder $a \bmod p$.
\end{frameddefn}

\begin{frameddefn}[label={def:field}]{Finite field $\mathbb{F}_p$}
  If $p$ is prime, the set $\mathbb{F}_p = \{0, 1, \ldots, p-1\}$ with addition and multiplication modulo $p$ is a \textbf{field}: we can add, subtract, multiply and \textbf{divide} by any non-zero element.
\end{frameddefn}

Being a field is what makes elliptic curves work: the formulas for point addition contain a \textit{division} (the slope of a line), and division is only always possible when the modulus is prime.

\subsection{Division = multiplication by the inverse}

Dividing by $d$ means multiplying by its \textbf{modular inverse} $d^{-1}$, i.e. the unique element such that $d \cdot d^{-1} \equiv 1 \pmod p$.

\begin{framedthm}[label={thm:fermat}]{Fermat's little theorem}
  If $p$ is prime and $d \not\equiv 0 \pmod p$, then $d^{p-1} \equiv 1 \pmod p$. Consequently
  \[
    d^{-1} \equiv d^{p-2} \pmod p .
  \]
\end{framedthm}

In code there are two standard ways to compute the inverse:
\begin{itemize}
  \item \textbf{Fermat}: $d^{p-2} \bmod p$ with fast modular exponentiation (simple, fine for a prime modulus);
  \item \textbf{Extended Euclidean algorithm}: finds $x, y$ with $d x + p y = \gcd(d,p) = 1$, and then $x \equiv d^{-1}$. This is the method normally used in practice.
\end{itemize}

\begin{framedex}{Inverses modulo 17}
  $10^{-1} \bmod 17 = 12$, because $10 \cdot 12 = 120 = 7 \cdot 17 + 1$.
  $\;(-4)^{-1} = 13^{-1} = 4$, because $13 \cdot 4 = 52 = 3 \cdot 17 + 1$.
  $\;(-1)^{-1} = 16^{-1} = 16$, because $16 \cdot 16 = 256 = 15 \cdot 17 + 1$.
\end{framedex}

\begin{framedobs}{A JavaScript pitfall}
  In JavaScript the operator \texttt{\%} is the \textit{remainder}, not the modulo: \texttt{-5 \% 17} gives \texttt{-5}, not \texttt{12}. A correct helper is
  \begin{lstlisting}[language=Java,numbers=none]
const mod = (a, p) => ((a % p) + p) % p;\end{lstlisting}
  For real-size keys (256 bits) we must use \texttt{BigInt} (e.g. \texttt{17n}), because ordinary \texttt{Number}s are exact only up to $2^{53}$.
\end{framedobs}

\subsection{Squares modulo $p$}

To find the points of a curve $y^2 = f(x)$ we need to know which values are perfect squares modulo $p$ (\textit{quadratic residues}). Modulo 17:
\begin{center}
  \begin{tabular}{c|cccccccc}
    \toprule
    $y$ & 1 & 2 & 3 & 4 & 5 & 6 & 7 & 8 \\
    $y^2 \bmod 17$ & 1 & 4 & 9 & 16 & 8 & 2 & 15 & 13 \\
    \bottomrule
  \end{tabular}
\end{center}
Since $(17 - y)^2 \equiv y^2$, the values $y = 9, \ldots, 16$ give the same squares. So the non-zero squares are $\{1, 2, 4, 8, 9, 13, 15, 16\}$: exactly $(p-1)/2 = 8$ of them, each with \textbf{two} square roots $\pm y$.
